%==========================================================================
% 基于动态图注意力网络与多智能体强化学习的低空交叉路口飞行器协同管控
% 期刊论文版式（模仿低空交通/智能交通领域权威期刊论文布局）
% 作者：彭义森  指导老师：韩云祥  四川大学
% 版本：v14（期刊版式）
% 编译方式：xelatex（两遍）
%==========================================================================
\documentclass[12pt,a4paper,UTF8]{ctexart}

%--------------------------------------------------------------------------
% 宏包
%--------------------------------------------------------------------------
\usepackage{geometry}
\usepackage{amsmath,amssymb,bm}
\usepackage{graphicx}
\usepackage{booktabs}
\usepackage{hyperref}
\usepackage{float}
\usepackage{enumitem}
\usepackage{multirow,array}
\usepackage{cite}
\usepackage{caption}
\usepackage{subcaption}
\usepackage{threeparttable}
\usepackage{algorithm}
\usepackage{algorithmic}
\usepackage{tikz}

\geometry{left=3cm,right=2.5cm,top=2.5cm,bottom=2.5cm}
\hypersetup{colorlinks=true,linkcolor=blue,citecolor=blue,urlcolor=blue}

\newcommand{\dyg}{DyGAT}

\begin{document}

%--------------------------------------------------------------------------
% 标题与作者
%--------------------------------------------------------------------------
\begin{center}
    {\LARGE\textbf{基于动态图注意力网络与多智能体强化学习的\\[2pt]低空交叉路口飞行器协同管控}}\\[10pt]
    {\large 彭义森，韩云祥}\\[4pt]
    {\small （四川大学，四川 成都 610065）}
\end{center}

%--------------------------------------------------------------------------
% 中文摘要（期刊式：单段，目的-方法-结果-结论）
%--------------------------------------------------------------------------
\vspace{8pt}
\noindent\textbf{摘\quad 要：}针对高密度、强动态低空空域中多架电动垂直起降飞行器（eVTOL）在交叉路口的协同安全通过难题，提出一种融合动态图注意力网络（DyGAT）的耦合多智能体强化学习协同管控模型。模型以分布式马尔可夫决策过程为形式化框架，核心创新是在经典图注意力网络中引入基于物理量的乘法门控注意力掩码：利用三个手工设计的指数核将飞行器对之间的距离、速度差与航向差编码为$(0,1]$区间的软性掩码信号，对学习注意力进行乘法门控，引导模型把有限的计算资源聚焦于真正构成威胁的交互对；同时设计速度自适应安全距离与分层奖励函数，并结合课程学习与Episode级冲突损失加权，缓解稀疏冲突样本导致的训练困难。在21架eVTOL的二维交叉路口仿真场景中，模型与8种基线方法进行了系统对比。实验结果表明：无冲突完成率（CFCR）从最优基线MAPPO的$57.8\%$提升至$68.4\%$（$p<0.001$，Cohen's $d=5.27$），最小间隔违反率（MSVR）从$3.0\%$降至$2.4\%$（$p=0.029$），均达到统计显著水平；消融实验表明DyGAT是贡献最大的模块，移除后总冲突次数增加$173.6\%$。研究结论表明，物理先验注入与学习注意力机制形成"物理划界、学习微调"的互补格局，为高密度低空交通的智能化协同管控提供了兼具安全性与可扩展性的技术路径。

\vspace{6pt}
\noindent\textbf{关键词：}低空交通管控；城市空中交通；多智能体强化学习；动态图注意力网络；乘法门控注意力掩码；冲突避免

%--------------------------------------------------------------------------
% 英文摘要
%--------------------------------------------------------------------------
\vspace{10pt}
\noindent\textbf{Abstract:}\enspace The safe and coordinated passage of multiple electric vertical take-off and landing (eVTOL) aircraft through urban low-altitude intersections is a critical challenge for high-density low-altitude airspace operations. This paper proposes a coupled multi-agent reinforcement learning (MARL) model enhanced by a dynamic graph attention network (DyGAT). The model is formalized within the decentralized Markov decision process framework, and its core innovation lies in a multiplicative gating attention mask driven purely by physical quantities: three hand-crafted exponential kernels encode the pairwise distance, speed difference, and heading difference into a soft mask in $(0,1]$, multiplicatively gating the learned attention so that limited representational capacity is concentrated on genuinely hazardous interaction pairs. A speed-adaptive safety distance, a hierarchical reward function, curriculum learning, and episode-level conflict loss weighting are further designed to improve learning efficiency under sparse conflict samples. In a simulation with 21 eVTOLs crossing a two-dimensional intersection, the proposed method is systematically compared with eight baselines. The conflict-free completion rate (CFCR) is improved from $57.8\%$ (MAPPO) to $68.4\%$ ($p<0.001$, Cohen's $d=5.27$), and the minimum separation violation rate (MSVR) drops from $3.0\%$ to $2.4\%$ ($p=0.029$); ablation studies show that removing DyGAT increases the total conflicts by $173.6\%$. These results suggest that injecting physics-structured priors into learned attention yields a complementary ``physics-delimiting, learning-refining'' paradigm, offering a safe and scalable pathway toward intelligent management of high-density low-altitude air traffic.

\vspace{6pt}
\noindent\textbf{Key words:}\enspace low-altitude traffic management; urban air mobility; multi-agent reinforcement learning; dynamic graph attention network; multiplicative gating attention mask; conflict avoidance

\vspace{12pt}
\hrule
\vspace{12pt}

%==========================================================================
% 1 引言
%==========================================================================
\section{引言}

\subsection{研究背景与意义}

近几年，低空经济被列入国家战略性新兴产业，电动垂直起降飞行器（Electric Vertical Take-Off and Landing，简称eVTOL）也随之从实验室原型快步走向商业示范。eVTOL是一种能像直升机一样垂直起降、又像固定翼飞机一样高效巡航的电动飞行器，不需要跑道，能在楼顶、停车场这类城市空间起降，因此被视为城市空中交通（Urban Air Mobility，UAM）\cite{ref45}的核心载运工具。亿航EH216-S已经在国内拿到了全球首张无人驾驶eVTOL型号合格证，Joby、Volocopter、Archer等国外企业也陆续进入适航审定阶段。美国联邦航空管理局（FAA）在《城市空中交通运行概念2.0》中预计，到2040年美国将有数百万架次的无人机和eVTOL在低空空域飞行\cite{ref1}。

机队规模一旦达到这种量级，一个现实问题就会立刻浮出水面：这么多飞行器挤在几百米以下的低空空域里，靠什么来保证它们互不碰撞、有序通行？传统民航空管靠管制员语音指挥，间隔以公里计，根本应对不了低空这种"一平方公里几十架、状态秒级变化"的运行密度。低空交通管理系统（UTM）因此被寄予厚望——它要求用自动化系统实时感知空域态势、自动做冲突探测与解脱。而在所有低空场景中，交叉路口又是冲突最集中、最棘手的地方：多股交通流在路口中心交汇，任何两架相邻方向飞来的飞行器，航向线都会在中心附近交叉，冲突几乎不可避免。如何让这些飞行器在保证安全间隔的前提下尽量不打乱节奏地通过路口，正是先进空中交通（Advanced Air Mobility，AAM）落地前必须迈过的一道坎。

从研究的角度看，这个问题本质上是"多架飞行器在共享空域里的协同序贯决策"：每一架飞行器的最佳动作，都取决于其他飞行器接下来会怎么飞，彼此之间是典型的博弈耦合关系。它既要解决"安全"（别撞上），又要兼顾"效率"（别堵死），还要在决策周期只有几百毫秒的限制下实时求解。这恰好是图神经网络擅长建模结构化空间关系、注意力机制擅长筛选关键信息、多智能体强化学习（MARL）擅长序贯决策的交叉地带，也是本文的工作切入点。

\subsection{本文主要贡献}

针对低空交叉路口的飞行器协同管控问题，本文的主要贡献概括如下：

（1）构建了面向低空交叉路口冲突规避的仿真环境，引入速度自适应安全距离与eVTOL物理运动约束，并纳入风扰与传感器噪声，使仿真条件尽量贴近真实低空运行。

（2）提出动态图注意力网络（DyGAT）。它在经典图注意力网络（GAT）基础上，首次引入基于物理量的乘法门控注意力掩码——用三个手工设计的指数核，把飞行器对之间的距离、速度差和航向差编码成一个$(0,1]$区间的软性掩码，乘法作用于学习注意力，引导模型聚焦高风险交互对。消融实验表明，去掉这一模块后总冲突次数增加$173.6\%$。

（3）设计课程学习与Episode级冲突损失加权相结合的训练机制，在不破坏PPO on-policy框架的前提下，显著缓解了稀疏冲突样本带来的训练困难。

（4）建立双层评价指标体系，揭示PCR（$99.1\%$）与CFCR（$68.4\%$）之间约30个百分点的系统性差距，并对CTDE-MARL安全评估方法论提出警示。

（5）与8种基线方法系统对比，给出包含Welch $t$检验、Cohen's $d$效应量与统计功效分析的完整统计证据链。

本文其余内容组织如下：第2节梳理相关工作并指出已有研究的不足；第3节将问题形式化为分布式马尔可夫决策过程并介绍仿真环境与奖励函数设计；第4节详细阐述DyGAT-MARL模型的结构与训练机制；第5节给出实验设置与结果分析；第6节进行讨论并总结全文。

%==========================================================================
% 2 相关工作
%==========================================================================
\section{相关工作}

围绕低空/空中交通的协同管控，学术界已经形成了"传统优化—工程安全—数据驱动"三条比较清晰的路线。本节沿着这条脉络梳理已有研究，并顺带把本文用到的基础概念解释清楚，为后文铺垫。

\subsection{传统空中交通流量管理方法}

空中交通流量管理（Air Traffic Flow Management，ATFM）是空管领域最早成体系的优化框架，核心思路是在已知空域容量的前提下，用整数规划、线性规划等数学手段，对航班的地面等待、改航和调速做集中式决策。Menon等\cite{ref2}把航空交通流抽象成连续介质模型，用偏微分方程描述交通流密度随时间的演化，再求解最优流量控制策略。这套方法在高空航路网络里很有效，但前提是"空域状态全局可知、决策由单一中心节点完成"。一旦飞行器数量上升到几十上百架，状态空间和约束规模就会指数膨胀，实时求解基本不可能。这种"集中式优化难以扩展"的固有短板，正是后来数据驱动方法切入的地方。

值得一提的是，地面交通信号控制领域的研究为空中交通管控提供了重要的方法论借鉴。作为自适应信号控制的代表工作，Aslani等\cite{ref12}把actor-critic强化学习算法用到真实交通网络的信号控制上，验证了它在多种交通扰动下的鲁棒性；国内学者也较早探索了基于Q-learning的交叉口信号控制\cite{ref16,ref19,ref20}，并逐步过渡到深度强化学习\cite{ref15,ref17,ref18,ref22}。这些工作表明，强化学习在处理多路口、多时段、强耦合的交通控制问题方面具有优势，这也为本文将其思想推广到空中交通场景提供了依据。

\subsection{基于工程约束的冲突解脱方法}

与宏观流量管理相对，工程安全方法着眼于微观层面的冲突解脱。它们的共同点是拿明确的物理或数学约束做安全保证，但各有各的适用边界。

\textbf{模型预测控制（MPC）。}Richards与How\cite{ref3}把飞行器航迹规划与冲突避免建模成混合整数线性规划，在滚动时域里反复求解最优控制序列。打个比方，MPC就像下棋时多想几步再落子：每走一步都基于当前局面预测未来一段时间的走向，选出最优方案，但只执行第一步，然后再重新计算。它的优点是约束表达清晰、最优性可度量，缺点是计算量随预测时域和飞行器数量快速增长，在高密度低空场景下实时性难以保证。

\textbf{速度障碍法（VO）与最优互惠避碰（ORCA）。}速度障碍法的思想很直观：如果两架飞行器保持当前相对速度不变就会相撞，那么所有"会导致相撞的相对速度"构成一个障碍集合，避碰就是选择一个不落在集合里的速度。van den Berg等\cite{ref4}在此基础上提出ORCA，引入"互惠"假设——双方各承担一半避让责任，然后通过线性规划选择偏离期望速度最小的安全速度，实现了分布式的平滑避碰。ORCA已广泛用于机器人与无人机导航，但它依赖"所有智能体对称地分担避让责任"这一假设，一旦智能体行为异构或目标相互冲突（例如一方正在避让、另一方却在争抢通行权），避让质量就会明显下降。

\textbf{控制屏障函数（CBF）。}Ames等\cite{ref5}系统总结了CBF理论。它的核心思想是设计一个满足前向不变性条件的"安全屏障"函数，把安全约束编码成在线可解的二次规划约束，在尽量不干扰名义控制器的前提下保证系统始终待在安全集合内。CBF的安全保证是严格的，但对系统模型精度要求很高，而且多架飞行器共享空域时，各个安全集合的交集会迅速收缩，可行解很容易退化。

总的来说，工程安全方法在形式化安全保证方面具有不可替代的优势，但面对高密度、强动态的交互场景都不同程度地存在局限，这为学习方法留出了空间。

\subsection{强化学习与多智能体强化学习在交通管控中的应用}

强化学习（Reinforcement Learning，RL）让智能体通过与环境的试错交互来学习策略，不需要精确建模系统动力学。它背后的数学模型是马尔可夫决策过程（MDP）：用状态、动作、转移概率、奖励和折扣因子五元组描述"智能体在什么状态、做什么动作、获得多少回报、环境怎么变"这一整套循环。智能体的目标，是学到一个能最大化长期累计回报的策略。

在单智能体层面，深度Q网络（DQN）\cite{ref40}及其变体被广泛用于冲突规避和信号控制\cite{ref12,ref15,ref16,ref17,ref18,ref19,ref20,ref21,ref22}。江波等\cite{ref21}将深度Q学习应用于机场道口的协同调速，验证了RL在航空地面场景的可行性。但单智能体RL存在一个固有局限：它把其他飞行器当作环境的一部分，无法刻画它们之间的博弈耦合关系。当每架飞行器都只优化自身收益时，联合行为往往偏离系统最优，甚至引发冲突的连锁反应。

多智能体强化学习（Multi-Agent Reinforcement Learning，MARL）正是为解决这个问题而生的。多智能体场景的核心困难在于非平稳性：其他智能体的策略也在不断更新，使得每个智能体面对的环境始终处于动态变化之中，破坏了单智能体强化学习所依赖的马尔可夫性。另一大困难是信用分配，即如何把团队共同获得的收益，合理归因到每个智能体的具体行为上。为应对这些难题，学术界发展出集中训练分布式执行（Centralized Training with Decentralized Execution，CTDE）范式：训练阶段共享全局信息，把智能体间的耦合关系显式建模，使训练信号更稳定；执行阶段每个智能体仅依据自身局部观测独立决策，以满足实时性与通信约束。围绕CTDE，Tumer与Agogino\cite{ref6,ref7,ref14,ref24}较早验证了分布式多智能体学习在ATM中的有效性；Spatharis等\cite{ref11}、Ghavamzadeh等\cite{ref23}设计了分层MARL方案；Brittain与Wei\cite{ref8}在高密度航路扇区验证了深度MARL的自主间隔保障；Ghosh等\cite{ref9}用深度集成MARL提升了管制决策的稳健性，Xie等\cite{ref29}则将强化学习用于UAM与密集低空交通的流量管理。算法层面，MADDPG（Multi-Agent Deep Deterministic Policy Gradient）\cite{ref34}采用集中式Critic与分布式Actor的结构，训练时Critic借助全局信息评估联合动作的价值，执行时每个Actor仅依据自身观测输出动作；COMA（Counterfactual Multi-Agent Policy Gradients）\cite{ref35}通过反事实基线——比较"智能体实际动作带来的回报"与"若它换作其他动作的期望回报"——衡量每个智能体对团队收益的真实边际贡献，从而解决信用分配问题；QMIX\cite{ref33}通过单调混合网络将全局Q值分解为个体Q值的单调组合，要求全局Q值随任一个体Q值单调不减，以保证分解与全局最优一致；MAPPO（Multi-Agent PPO）\cite{ref36}则把单智能体PPO的稳定训练特性推广到多智能体场景，是目前合作博弈任务中的强基线。国内方面，翟子洋等\cite{ref13}系统梳理了大规模智慧交通信号控制中强化学习的技术演进。

\subsection{图神经网络与注意力机制}

飞行器之间的关系天然适合用"图"来表达：每架飞行器是节点，飞行器对之间的交互是边。图神经网络（GNN）就是专门在这种结构上做特征学习的模型。Kipf与Welling\cite{ref32}提出的图卷积网络（GCN）通过邻居聚合更新节点特征，但聚合权重只由图的拓扑结构（即节点度数，也就是邻居数量）决定，不随节点特征变化，因此难以区分重要邻居与次要邻居。Veličković等\cite{ref31}提出的图注意力网络（GAT）解决了这个问题——它让网络自己学习"哪些邻居更值得关注"，通过注意力权重自适应地聚合信息。

注意力机制背后还有一个与本文密切相关的重要思想：门控。长短期记忆网络（LSTM）\cite{ref37}之所以能记住长序列信息，靠的就是遗忘门、输入门、输出门这一套门控结构，通过控制"历史信息保留多少、新信息写入多少"来调节信息流。这类门控信号都是通过训练学出来的。本文尝试一个相反的思路：不用学习，直接让物理量充当门控信号，把飞行中的基本安全常识注入注意力计算，这正是后文DyGAT的核心思想。

在低空无人机交通领域，Li等\cite{ref10}提出的MS-GRL方法用多尺度GAT增强强化学习，解决了100架UAV的密集冲突消解问题，是图增强MARL的代表性工作。但它和大多数GNN方法一样，图结构基于固定距离阈值构建——只要两机距离在阈值以内就一律视为邻居。其局限在于：它区分不了"距离相同但风险不同"的飞行器对。两对飞行器当前欧氏距离完全一样，一对相向高速靠近、另一对同向等速，冲突风险能差出几十倍，固定阈值图却把它们一视同仁。这正是本文DyGAT要解决的问题。

\subsection{现有研究的不足与本文切入点}

综合上面几节，已有研究在三个维度上还存在明显缺口：

（1）\textbf{时空耦合建模不足。}多数方法用全连接网络或静态图结构\cite{ref8,ref10,ref11}，区分不了"距离相同但相对运动状态不同"的飞行器对，在高密度动态场景下决策质量受限。

（2）\textbf{稀疏冲突样本利用不足。}冲突是低概率事件，在训练数据里占比通常不到$0.1\%$\cite{ref12,ref13}，随机采样下这些宝贵的危险样本很容易被大量正常样本淹没，难以提升模型对罕见危险的敏感性。

（3）\textbf{环境不确定性建模不足。}多数研究忽略物理运动约束和动态安全距离\cite{ref11}，风扰和噪声处理得也过于简单。安全间隔如果是个固定常数，就反映不出"飞得越快、越要多留余地"这个基本规律。

正是这三个缺口，构成了本文工作的出发点。

%==========================================================================
% 3 问题建模与管控模型设计
%==========================================================================
\section{问题建模与管控模型设计}

\subsection{问题描述与场景设定}

本文聚焦城市低空十字交叉路口这一典型高冲突场景。如图\ref{fig:scenario}所示，$N$架eVTOL均匀分布在以交叉中心为圆心、半径$r_{\text{init}}$的圆周上，各自从不同方位径直飞向中心。飞行器的航向线天然在中心区域附近相互交叉——任何两架相邻或相对方位出发的飞行器都会在路口附近"碰面"，冲突威胁持续存在。飞行器需要在保持安全间隔的前提下，以尽可能小的延迟通过交叉区域。

\begin{figure}[H]
    \centering
    \begin{minipage}[c]{0.7\textwidth}
    \centering
    \begin{tikzpicture}[scale=0.95]
        \draw[->,thick] (-5.2,0) -- (5.2,0);
        \draw[->,thick] (0,-5.2) -- (0,5.2);
        \draw[dashed] circle (3.2);
        \foreach \ang in {0,60,120,180,240,300} {
            \fill[blue!70] ({3.4*cos(\ang)},{3.4*sin(\ang)}) circle (0.12);
        }
        \foreach \ang in {0,60,120,180,240,300} {
            \draw[->,blue!70,thick] ({3.2*cos(\ang)},{3.2*sin(\ang)}) -- ({1.1*cos(\ang)},{1.1*sin(\ang)});
        }
        \draw[red,very thick] (1.0,0.35) circle (0.6);
        \node[below right] at (1.7,1.05) {冲突区};
        \node[above] at (0,5.3) {交叉路口中心};
        \node[below] at (0,-5.7) {初始位置圆周};
    \end{tikzpicture}
    \end{minipage}
    \caption{低空十字交叉路口场景示意图。蓝色圆点为均匀分布于圆周上的eVTOL，箭头指向交叉中心，红色圆环为高冲突区域。}
    \label{fig:scenario}
\end{figure}

这个场景虽然抽象，但抓住了低空交通冲突的三个本质特征。第一是\textbf{对称冲突}：所有飞行器以相同优先级争夺中心区域，不存在天然的"让行次序"。第二是\textbf{密集交互}：$N$架飞行器最多构成$N(N-1)/2$个交互对，$N=21$时就是210对，远超人工调度的能力。第三是\textbf{时空耦合}：冲突风险不仅看当前距离，还看相对速度和相对航向——同样是相距100 m，相向而行和同向而行的风险天差地别。

为保证问题可解，本文作如下合理假设：所有飞行器在同一高度层运行（二维平面问题）；各eVTOL具有相同的动力学参数（同构智能体）；每架飞行器只能感知500 m半径内的邻居（局部可观测性），全局状态仅在训练阶段可用；训练阶段通信无延迟、无丢包。这些假设与真实eVTOL的机载感知能力基本一致，假设的放宽讨论见第6.3节。

\subsection{分布式马尔可夫决策过程形式化}

把上述问题写成分布式马尔可夫决策过程（Dec-MDP）\cite{ref41}六元组$\langle \mathcal{N}, \mathcal{S}, \mathcal{A}, \mathcal{P}, \mathcal{R}, \gamma \rangle$：

\begin{itemize}
    \item \textbf{智能体集合} $\mathcal{N} = \{1, 2, \ldots, N\}$。训练阶段$N$由课程学习从5逐步增至21，评估阶段固定为21。
    \item \textbf{全局状态} $\mathcal{S}$：$S = (s_1, \ldots, s_N)$，其中$s_i = [x_i, y_i, v_i, v_i^{\text{target}}, d_i^{\text{center}}]$为飞行器$i$的状态向量，含位置、当前速度、目标速度与到中心的距离。需要说明的是，这个五维向量只是形式化层面的定义；策略网络实际输入的是扩展后的局部观测$o_i$（第3.3.1节）。全局状态$S$只在CTDE训练阶段供GAT-Mixer计算全局价值函数使用。
    \item \textbf{动作} $\mathcal{A}$：每架飞行器的动作为连续加速度$a_i$（m/s$^2$），受运动学约束、风扰与噪声共同影响。
    \item \textbf{转移概率} $\mathcal{P}$：由第3.3节的运动学方程与不确定性模型隐式确定，不显式建模。
    \item \textbf{奖励函数} $\mathcal{R}$：团队奖励结构$R = \frac{1}{N}\sum_i r_i$，另有完成奖励（第3.4节）。
    \item \textbf{折扣因子} $\gamma = 0.99$，用于权衡近期奖励与远期奖励的相对重要程度，$\gamma$越接近1，模型越重视长期收益。
\end{itemize}

这里要特别解释一下为什么用团队平均奖励。如果每架飞行器只优化自己的安全，很容易出现"把风险甩给邻居"的机会主义——我减速让你撞上，只要我不撞就行。把个体奖励平均成团队奖励后，每一架的行为都会影响全局收益，客观上激励飞行器采取有利于整体交通流的策略。

\subsection{仿真环境设计}

仿真环境（AirTrafficEnv）以Python面向对象方式实现，是全部实验的基础平台。下面按观测空间、动作空间、不确定性建模三部分展开。

\subsubsection{观测空间设计}

观测空间分三个层次，层层递进。第一层是形式化状态$s_i\in\mathbb{R}^5$，即上面定义的五维向量。第二层是局部观测$o_i$，在$s_i$基础上加入两路信息：空间上，加入500 m半径内各邻居$j$的相对位置、速度差、加速度与航向角，用于刻画邻域内的交互态势；时间上，加入自身最近$T=10$步（对应5 s）的历史状态序列，历史信息让模型能感知自己的运动趋势，这是判断未来风险的基础。第三层是编码特征，局部观测先经64维LSTM编码轨迹，再经256维DyGAT聚合空间信息，得到蕴含时空耦合信息的256维特征$h_i^{\text{coupled}}$：
\begin{equation}
    o_i \xrightarrow{\text{LSTM}_{64}} h_i^{\text{traj}} \xrightarrow{\text{DyGAT}_{256}} h_i^{\text{coupled}} \rightarrow \text{策略/Q值输出}.
    \label{eq:encoding_chain}
\end{equation}

\subsubsection{动作空间与物理运动约束}

动作空间为连续加速度$a_i\in[-4.0, 3.0]$ m/s$^2$，正值加速、负值减速。速度更新式为：
\begin{equation}
    v_i^{t+1} = \text{clip}\!\left(v_i^{t} + a_i \cdot \Delta t \cdot 3.6,\; 60,\; 140\right),
    \label{eq:velocity_update}
\end{equation}
其中$\Delta t=0.5$ s为控制周期，系数3.6用于把m/s换算成km/h，速度限制在60--140 km/h。此外，单步航向变化不超过$\pi/6$（30°），防止飞行器"瞬间掉头"这类不真实的运动。速度范围参考了当前主流eVTOL的巡航速度；加速度取3.0、减速度取4.0 m/s$^2$，减速度大于加速度是安全设计惯例。这里需要诚实说明：主流eVTOL厂商（Joby、Volocopter、Lilium、Archer等）都没有公开最大水平加速度的精确规格，学术界建模一般假设约2.0 m/s$^2$\cite{ref25}，本文取值是在缺乏精确数据下的工程估计，拿到厂商数据后应做敏感性验证。

\subsubsection{不确定性建模与动态安全距离}

真实低空环境存在两类不可回避的不确定性：传感器噪声与风扰。本文在观测通道叠加$\mathcal{N}(0, 0.5^2)$ m的高斯噪声，模拟GPS/视觉定位的典型误差；每架飞行器每步受到方向均匀随机、强度0.1 m/s的风扰。需要提前说明一个已知局限：真实低空风场通常用Dryden（MIL-F-8785C）\cite{ref26}或Kaimal（IEC 61400-1）\cite{ref44}谱模型描述，具有空间相关性和时间连续性，本文采用的均匀随机风扰是简化版本，这一局限将在第6.3节进一步讨论。

安全间隔同样不该是个固定常数。直观上讲，飞得越快，单位时间内接近的距离越大，需要留的余量就该越多。据此定义速度自适应安全距离：
\begin{equation}
    d_{\text{safety}} = d_{\text{base}} + \bar{v} \times 0.5,
    \label{eq:safety_distance}
\end{equation}
其中$d_{\text{base}}=50$ m为基本安全间隔，$\bar{v}$为当前时刻场景内飞行器的平均速度，$\bar{v}\times0.5$可理解为"额外0.5 s前瞻时间内扫过的距离"。以$\bar{v}=70$ km/h（约19.4 m/s）为例，$d_{\text{safety}}\approx60$ m；加速到140 km/h时约69 m。低速不保守、高速有保护，这与空管实践中"速度越高、间隔越大"的原则完全一致。

\subsection{分层奖励函数设计}

奖励函数是强化学习中连接行为与目标的桥梁，其设计质量直接影响学习效果。本文遵循三条原则：奖励信号要密集，以保证模型能频繁获得学习反馈；惩罚梯度要符合物理直觉，即越接近危险越重；多个目标之间要有明确的权重关系，使安全、效率与协同可以权衡。据此构造四部分奖励：
\begin{equation}
    r_i = r_{\text{safe}} + r_{\text{speed}} + r_{\text{coop}} + r_{\text{eff}}.
    \label{eq:reward_total}
\end{equation}

\textbf{安全奖励}采用三段式分级结构：
\begin{equation}
    r_{\text{safe}} = \begin{cases}
        -50 \exp\!\left(2\frac{d_{\text{safety}}-d_{\min}}{d_{\text{safety}}}\right), & d_{\min}<d_{\text{safety}} \text{（危险区）}\\
        -10 \frac{1.2d_{\text{safety}}-d_{\min}}{0.2d_{\text{safety}}}, & d_{\text{safety}}\leq d_{\min}<1.2d_{\text{safety}} \text{（缓冲区）}\\
        +5, & d_{\min}\geq1.2d_{\text{safety}} \text{（安全区）}
    \end{cases}
    \label{eq:reward_safe}
\end{equation}
其中$d_{\min}$为飞行器$i$到所有邻居的最小距离。三个区域的物理含义很清晰：已经侵入安全间隔（危险区）时采用指数惩罚，距离越近惩罚越重，引导模型尽快拉开距离；处于警示缓冲带时采用线性惩罚，鼓励多留余量；保持充足间隔则给予正奖励。三段式的设计避免了单一线性惩罚在安全边界处突变、梯度不连续的问题。

\textbf{速度奖励}引导飞行器把速度维持在理想区间$[65,90]$ km/h（与目标速度$70\pm10$ km/h一致）：
\begin{equation}
    r_{\text{speed}} = \begin{cases}
        -50\frac{v_i-90}{50}, & v_i>90 \text{（超速）}\\
        -30\frac{65-v_i}{25}, & v_i<65 \text{（过慢）}\\
        +20, & 65\leq v_i\leq 90 \text{（理想）}
    \end{cases}
    \label{eq:reward_speed}
\end{equation}
超速会放大冲突风险，过慢会拖累交通流，只有落在理想区间才能兼顾安全与效率。

\textbf{协同奖励}刻画与邻居的速度一致性：
\begin{equation}
    r_{\text{coop}} = -0.2\left|v_i - \bar{v}_{\text{neighbor}}\right|,
    \label{eq:reward_coop}
\end{equation}
其中$\bar{v}_{\text{neighbor}}$为邻域内邻居的平均速度。这项奖励意在抑制"个别飞行器突进、整体节奏被打乱"的现象，在路口场景里有助于形成平滑的交通流。

\textbf{效率奖励}引导飞行器逼近各自的目标速度：
\begin{equation}
    r_{\text{eff}} = -0.5\left|v_i - v_i^{\text{target}}\right|,
    \label{eq:reward_eff}
\end{equation}
保证飞行器在规避冲突的同时仍以接近任务速度的方式飞行，避免"为了安全而过度保守"。

全局奖励为个体奖励的平均，$R=\frac{1}{N}\sum_i r_i$；当所有飞行器都通过交叉中心（距中心小于10 m）时，额外奖励+200。这个完成奖励属于稀疏奖励——只在Episode结束时出现一次——但幅度足够大，能让模型把"完成任务"这一长期目标纳入优化考量，与高密度的即时奖励形成互补。

%==========================================================================
% 4 基于DyGAT-MARL的协同管控模型
%==========================================================================
\section{基于DyGAT-MARL的协同管控模型}

\subsection{模型总体架构}

模型采用集中训练分布式执行（CTDE）范式，总体架构如图\ref{fig:architecture}所示，从下到上分四层：环境层（AirTrafficEnv，含风扰与噪声）、编码层（LSTM时序编码+ DyGAT空间聚合，含乘法门控注意力掩码）、决策层（策略头+ Q值头+GAT-Mixer，训练集中、执行分布）、训练层（课程学习+冲突损失加权+PPO）。

\begin{figure}[H]
    \centering
    \includegraphics[width=1.0\textwidth]{QQ图片20260523201329.png}
    \caption{模型总体架构——环境层、编码层（LSTM+DyGAT含乘法门控注意力掩码）、决策层（策略头+Q值头+GAT-Mixer）与训练层（课程学习+冲突损失加权+PPO），遵循集中训练分布式执行范式。}
    \label{fig:architecture}
\end{figure}

四层架构的设计思想可以概括成一句话：环境层保真，编码层求精，决策层分布，训练层集中。环境层尽量贴近真实运行；编码层通过时空注意力与物理门控，从观测里提炼决策所需的关键信息；决策层保证执行阶段的实时性与分布式特性；训练层则靠课程式学习与冲突加权，化解稀疏奖励下的优化难题。

\subsection{LSTM时序编码}

飞行器的运动有明显的时序连贯性，当前位置和速度的历史轨迹里蕴含着运动趋势，是判断未来风险的重要线索。模型用两层LSTM（隐层64维，dropout 0.1）做时序编码：
\begin{equation}
    h_i^{\text{traj}} = \text{LSTM}_{64}(o_i^{t-T+1:t}),
    \label{eq:lstm}
\end{equation}
其中$o_i^{t-T+1:t}$为最近$T=10$步的观测序列，$h_i^{\text{traj}}\in\mathbb{R}^{64}$为轨迹嵌入。LSTM的记忆单元使模型能够保留飞行器速度变化趋势与转向意图等时序信息——例如，一个正在减速避让的邻居与一个即将加速穿越的邻居，二者接下来的行为截然不同，仅凭当前瞬时状态无法区分，需要借助历史轨迹才能判断。

\subsection{动态图注意力网络（DyGAT）}

DyGAT是本文的核心创新模块。它输入$N$个节点的特征（以LSTM输出$h_i^{\text{traj}}\in\mathbb{R}^{64}$为节点特征，输出维度$F_{\text{out}}=256$），输出融合了空间与时序信息、并经物理门控调制后的节点表示。模块由空间注意力、时序注意力和乘法门控注意力掩码三个组件构成。

\subsubsection{空间注意力}

空间注意力沿用GAT的标准范式，衡量邻居节点对当前节点信息的重要性。节点$i$与节点$j$之间的注意力分数为（其中LeakyReLU$_{0.2}$表示负半轴斜率为0.2的带泄漏修正线性单元）：
\begin{equation}
    e_{ij} = \text{LeakyReLU}_{0.2}\!\left(a^{\top}[W h_i^{\text{traj}} \parallel W h_j^{\text{traj}}]\right),
    \label{eq:spatial_score}
\end{equation}
其中$W$为共享线性变换，$\parallel$为向量拼接，$a$为注意力参数向量。归一化后得到注意力权重：
\begin{equation}
    \alpha_{ij} = \frac{\exp(e_{ij})}{\sum_{k\in\mathcal{N}_i}\exp(e_{ik})},
    \label{eq:spatial_softmax}
\end{equation}
节点$i$的空间聚合特征为：
\begin{equation}
    h_i^{\text{spatial}} = \sum_{j\in\mathcal{N}_i}\alpha_{ij}\, W h_j^{\text{traj}}.
    \label{eq:spatial_agg}
\end{equation}
空间注意力实现了由数据自适应决定"哪些邻居更值得关注"的聚合方式，但正如第2.4节所述，纯数据驱动的注意力在训练数据稀疏时可能偏离物理直觉——这正为后文的乘法门控机制提供了动机。

\subsubsection{时序注意力}

为利用多步历史信息，DyGAT在注意力头上引入时序聚合。对每架飞行器，把最近$T=10$步的空间编码特征按匹配度加权求和：
\begin{equation}
    h_i^{\text{temporal}} = \sum_{\tau=1}^{T} \beta_{i,\tau}\, h_i^{\text{spatial}}(t-\tau+1),
    \label{eq:temporal_attn}
\end{equation}
其中$\beta_{i,\tau}$为时序注意力权重，由当前特征与第$\tau$个历史时刻特征的相似程度经softmax归一化得到，数值越大表示该历史时刻对当前决策的参考价值越高。时空融合表示为：
\begin{equation}
    h'_i = h_i^{\text{spatial}} + \alpha_{\text{temp}} \cdot h_i^{\text{temporal}},
    \label{eq:temporal_fusion}
\end{equation}
融合系数$\alpha_{\text{temp}}=0.3$为经验值，说明当前空间信息占主导、历史信息起辅助修正作用——这也符合"决策主要依据当前态势"的直觉。

\subsubsection{乘法门控注意力掩码（核心创新）}

上述空间与时序注意力都是纯学习机制。纯学习机制的隐患在于：当训练数据没覆盖所有交互模式时，学习注意力可能产生"反物理"的分配——比如给远距离、航向正交的飞行器很高的注意力。为从机制上约束这种风险，DyGAT用三个手工设计的指数核，把飞行器对之间的物理量编码成乘法门控信号$w_{ij}$：
\begin{equation}
    w_{ij} = \exp\!\left(-\frac{d_{ij}}{1000}\right) \cdot \exp\!\left(-\frac{|v_i - v_j|}{20}\right) \cdot \exp\!\left(-\frac{|\theta_i - \theta_j|}{\pi/4}\right),
    \label{eq:gating_w}
\end{equation}
并以乘法方式作用于学习注意力：
\begin{equation}
    e_{ij}^{\text{final}} = e_{ij} \odot w_{ij}.
    \label{eq:gating_final}
\end{equation}

$w_{ij}$的三个因子分别对应三条基本的安全常识：距离越近，威胁越大（距离因子）；速度差越大，相对运动越剧烈，碰撞风险越高（速度差因子）；航向差越大，越接近迎面交叉，冲突可能性越大（航向差因子）。三者相乘，构成完整的门控信号。

为直观说明门控效果，表\ref{tab:gating}给出四种典型飞行器对情景下的$w_{ij}$数值（$N=21$，$d_{\text{safety}}\approx85$ m）。

\begin{table}[H]
    \centering\small
    \caption{$w_{ij}$在不同飞行器对情景下的门控效果}
    \label{tab:gating}
    \begin{tabular}{|c|c|c|c|c|c|}
        \hline
        \textbf{情景} & $d_{ij}$(m) & $|v_i-v_j|$(km/h) & $|\theta_i-\theta_j|$ & $w_{ij}$ & \textbf{门控效果} \\
        \hline
        近距离同向等速 & 100 & 5 & $\pi/12$ & $\approx0.85$ & 轻度衰减（17\%），$e_{ij}$可自由调节 \\
        中距离同向等速 & 300 & 10 & $\pi/8$ & $\approx0.58$ & 中度衰减（42\%） \\
        中距离相向高速 & 200 & 40 & $\pi/2$ & $\approx0.015$ & 强衰减（98.5\%），注意力降至1/67 \\
        远距离航向正交 & 500 & 30 & $\pi/2$ & $\approx0.003$ & 近零衰减（99.7\%），信息流几乎阻断 \\
        \hline
    \end{tabular}
\end{table}

表中的对比极具说服力：两对飞行器距离同为200 m左右，同向等速时$w_{ij}\approx0.85$，注意力几乎不受限制；相向高速时$w_{ij}\approx0.015$，注意力被压到1/67。而固定阈值GAT的邻接矩阵里，二者同为"邻居"，权重完全一样——这正是DyGAT相对静态图方法的本质优势。

\subsubsection{门控机制的性质分析}

乘法门控有三个重要性质，共同保证了模型的物理合理性。

\textbf{软性非阻断。}$w_{ij}\in(0,1]$，物理先验不会完全掐断信息流。即使最不利的情形（远距离+高速差+航向正交），$w_{ij}$也只是一个极小正数而非零，给学习注意力保留了在极端条件下发现非常规依赖的可能。硬性掩码"一票否决"导致的表达力损失在这里不会发生。

\textbf{非对称性。}乘法形式（$\odot$而非$+$）使门控具有非对称性：$w_{ij}$可以近乎完全压制$e_{ij}$，但$e_{ij}$永远没法"抬升"一个已经被压到接近零的权重。换句话说，物理安全约束是"学习覆盖不了"的——就算训练数据里出现反例，模型也没法通过学习推翻物理门控划定的危险边界。这一点赋予了模型安全保证的鲁棒性。

\textbf{零参数零开销。}$w_{ij}$由状态量直接计算，不需要梯度回传。在$N=21$时，单步只增加约0.02 ms的额外开销，相对于约4.5 ms的整体推理时间几乎可以忽略，实时性完全不受影响。

从理论上看，这个门控机制本质上是\emph{结构化先验注入}（structured prior injection）。它和LSTM里的遗忘门$\sigma(Wx+b)$那种"学出来的门控"有本质区别：标准门控的门控信号是数据驱动的，本文的$w_{ij}$完全由物理量驱动、零可学习参数。更直观地说，$w_{ij}$划定了学习注意力可以操作的"有效范围"，学习注意力在这个范围内做精细化分配——二者形成"物理划界、学习微调"的互补格局。

\subsection{决策网络与GAT-Mixer混合网络}

每架飞行器维护一个单智能体网络，信息流为：LSTM编码器（2层，64维）$\to$ DyGAT层（输出256维$h_i^{\text{coupled}}$）$\to$ 策略头与Q值头。策略头是2层MLP（Tanh激活），输出动作均值$\mu_i$，动作按高斯分布$\mathcal{N}(\mu_i, 0.1^2)$采样；固定标准差0.1让策略具备可调节的探索性，又避免了方差学习的不稳定。Q值头是3层MLP（ReLU激活），输出个体Q值。策略头与Q值头共享编码层参数，好处有二：一是特征提取层参数量大幅减少，降低过拟合风险；二是策略与价值共享语义空间，有利于价值函数对策略改进的指导。

集中训练阶段，需要把各智能体信息融合成一个全局Q值$Q_{\text{tot}}$来评估联合状态-动作的质量。本文设计GAT-Mixer完成这一融合：
\begin{equation}
    Q_{\text{tot}} = \text{MLP}\!\left(\text{MeanPool}\!\left(\text{GAT}\left(\{s_i\}_{i=1}^{N}\right)\right)\right),
    \label{eq:gat_mixer}
\end{equation}
其中$\{s_i\}$为各智能体的4维状态向量，在全连接图上经GAT聚合后取均值池化，再经MLP输出。GAT-Mixer的设计动机有三：一是以GAT注意力替代QMIX的超网络，突破单调性约束。QMIX要求全局Q值随任一个体Q值单调不减（$\partial Q_{\text{tot}}/\partial Q_i\geq0$），这一约束简化了分解问题，但也限制了表达能力——例如，当某个智能体的个体Q值偏高、需要通过全局协调压低其权重时，单调性约束无法表达这种"个别压降"的调整。二是GAT注意力天然适应$N$的动态变化，契合课程学习中$N$从5到21的增长。三是均值池化具有平移不变性，使$Q_{\text{tot}}$对智能体排列顺序不敏感。需要澄清一点：本文的GAT-Mixer与ICLR 2023的GraphMixer\cite{ref28}只是名字相近，后者是基于MLP-Mixer的时序链路预测模型，与本文的全局价值混合功能无关。

\subsection{训练机制}

\subsubsection{课程学习}

如果一开始就让21架飞行器在紧凑空域里同时学习，模型面对的奖励信号极度稀疏——绝大多数Episode中根本不发生冲突，安全奖励的梯度贡献趋近于零，模型很难从中学会如何规避危险。课程学习（Curriculum Learning）\cite{ref39}的思路正是针对这一难题：先让模型在较简单的任务上掌握基础避让能力，再逐步过渡到困难任务，从而降低初始学习的难度。

本文的课程设计如下：从$N=5$架、$d_{\text{base}}=80$ m开始（宽松场景），连续3个周期满足"冲突率低于阈值且完成率高于80\%"就触发进阶：
\begin{equation}
    N \leftarrow \min(N+2, 21), \qquad d_{\text{base}} \leftarrow \max(d_{\text{base}}-5, 50).
    \label{eq:curriculum}
\end{equation}
每次进阶增加2架飞行器、缩小5 m基本安全间隔，共9个级别。为防训练中性能回退，设有回退机制：连续5个周期不满足完成率要求时，$N$回退1架、$d_{\text{base}}$回升5 m。前期在$N=5$时学会的基础避让技能（减速让行、保持间隔）在后期$N=21$场景中具有迁移性，显著缩短了整体训练时间。

\subsubsection{Episode级冲突损失加权}

即使有了课程学习，冲突仍是低概率事件。本文进一步引入Episode级冲突损失加权：对包含冲突的Episode施加更大的损失权重，使模型在参数更新时给予危险样本更高的权重。

在PPO\cite{ref30} on-policy框架下，总损失为：
\begin{equation}
    L_{\text{total}} = \lambda\left(L_{\text{policy}} + L_{q} + 0.5\,L_{\text{value}}\right),
    \label{eq:total_loss}
\end{equation}
其中$L_{\text{policy}}$为PPO裁剪策略损失，系数为：
\begin{equation}
    \lambda = \begin{cases}
        2.0, & \text{含冲突的Episode}\\
        1.0, & \text{无冲突的Episode}
    \end{cases}.
    \label{eq:lambda}
\end{equation}

这里需要澄清一个容易混淆的点：该机制与优先经验回放（Prioritized Experience Replay，PER）有本质区别。PER改变的是采样分布——高优先级样本被更频繁地采样，使采样分布偏离原始分布，需要用重要性采样校正；而本文的$\lambda$只是缩放损失幅度（$\nabla_\theta(\lambda L)=\lambda\nabla_\theta L$），不涉及对数据分布的改变，因此\textbf{严格保持了PPO的on-policy特性}——所有样本仍来自当前策略，只是含冲突样本的梯度贡献被放大了2倍。

另一个诚实的说明：$\lambda$与PPO的clip有交互。裁剪作用于$r_t(\theta)$层面，当某条transition的$r_t(\theta)$被clip截断时，$\lambda$对它就不再有效——被裁剪的目标值已经固定了。因此$\lambda$加权在训练早期（策略不稳定、clip截断率高）的实际效果可能弱于预期。PTR-PPO\cite{ref27}用截断重要性采样处理off-policy轨迹，是更系统的方案，可作为后续改进方向。

$\lambda$的取值直接影响训练效果，本文在$\{1.0, 1.5, 2.0, 2.5, 3.0\}$上做了实验（$n=5$），结果见表\ref{tab:lambda}。$\lambda=2.0$时TC最低（18.2）、CFCR最高（68.4\%），为最优；$\lambda=3.0$收敛最快（约300轮）但最终性能反而下降（CFCR降至57.3\%），说明权重过大导致优化振荡。当然，五个离散值的最优未必是连续空间的全局最优，更精细的调参可能带来进一步提升，但$\lambda=2.0$已足以支撑本文结论。

\begin{table}[H]
    \centering
    \caption{$\lambda$取值实验（$n=5$）}
    \label{tab:lambda}
    \begin{tabular}{|c|c|c|c|}
        \hline
        $\lambda$ & TC$\downarrow$ & CFCR(\%)$\uparrow$ & 收敛轮次 \\
        \hline
        1.0 & $24.6\pm2.8$ & $53.2\pm2.3$ & $\sim800$ \\
        1.5 & $21.8\pm2.5$ & $58.7\pm2.0$ & $\sim600$ \\
        \textbf{2.0} & $\mathbf{18.2\pm2.1}$ & $\mathbf{68.4\pm1.8}$ & $\sim400$ \\
        2.5 & $19.5\pm2.3$ & $64.1\pm1.9$ & $\sim350$ \\
        3.0 & $22.3\pm3.1$ & $57.3\pm2.5$ & $\sim300$ \\
        \hline
    \end{tabular}
\end{table}

训练中的优势函数采用广义优势估计（Generalized Advantage Estimation，GAE）\cite{ref38}计算，它通过对多个时间尺度上时序差分误差的加权平均估计"在给定状态下采取某动作相对平均水平的好坏程度"，在偏差与方差之间取得平衡（本文取$\gamma=0.99$、$\lambda_{\text{GAE}}=0.95$）。综合全部设计，完整训练流程如算法\ref{alg:training}所示。

\begin{algorithm}[H]
\caption{基于DyGAT-MARL的协同管控模型训练流程}
\label{alg:training}
\begin{algorithmic}[1]
\STATE 初始化课程级别$c=1$，$N=5$，$d_{\text{base}}=80$ m，策略参数$\theta$，学习率$\eta=5\times10^{-4}$
\FOR{episode $k=1,2,\ldots,K$}
    \STATE 初始化环境（$N$架eVTOL），重置智能体隐藏状态
    \STATE $ep\_data \leftarrow \emptyset$，$has\_conflict \leftarrow \text{False}$
    \REPEAT
        \STATE 各智能体获取局部观测$o_i$，经LSTM+DyGAT编码得$h_i^{\text{coupled}}$
        \STATE 策略头输出动作均值$\mu_i$，采样动作$a_i\sim\mathcal{N}(\mu_i,0.1^2)$
        \STATE 环境执行联合动作$(a_1,\ldots,a_N)$，返回奖励$R$与冲突信息
        \IF{发生冲突} \STATE $has\_conflict \leftarrow \text{True}$ \ENDIF
        \STATE 将$(o,a,\log\pi,\hat{A},R)$存入$ep\_data$
    \UNTIL{Episode结束}
    \STATE 计算GAE\cite{ref38}优势估计（$\gamma=0.99$，$\lambda_{\text{GAE}}=0.95$）
    \STATE 计算$\lambda = 2.0 \text{ if } has\_conflict \text{ else } 1.0$
    \STATE 以$\lambda$加权总损失（式(\ref{eq:total_loss})），更新策略参数$\theta$
    \STATE 课程学习判定：连续3周期满足进阶条件则$N\leftarrow\min(N+2,21)$，$d_{\text{base}}\leftarrow\max(d_{\text{base}}-5,50)$
\ENDFOR
\STATE \RETURN 训练后的策略参数$\theta$
\end{algorithmic}
\end{algorithm}

%==========================================================================
% 5 实验设计与结果分析
%==========================================================================
\section{实验设计与结果分析}

本节对模型进行全面实验验证。先交代实验设置，然后从训练收敛性、综合性能、安全指标、统计显著性、消融实验、密度鲁棒性、参数敏感性与推理效率八个维度逐项分析。

\subsection{实验设置}

\subsubsection{仿真环境与参数}

实验在自研AirTrafficEnv环境中进行，全部代码基于Python 3.11与PyTorch实现，运行于NVIDIA RTX 4060（8 GB）GPU与Intel i9-14900K CPU平台。核心参数见表\ref{tab:env}。其中空域1000 m×1000 m对应典型城市街区尺度，飞行器初始位置位于600 m半径的圆周上，留有充分机动空间；最大步长200步（100 s）保证即使多轮避让也有足够时间完成穿越。

\begin{table}[H]
    \centering\small
    \caption{仿真环境参数}
    \label{tab:env}
    \begin{tabular}{|c|c|c|c|}
        \hline
        \textbf{参数} & \textbf{取值} & \textbf{参数} & \textbf{取值} \\
        \hline
        空域范围 & $1000\text{m}\times1000\text{m}$ & 最大转弯角 & $\pi/6$ rad \\
        $d_{\text{base}}$（基本安全间隔） & 50 m & 时序窗口 $T$ & 10步（5 s） \\
        $\Delta t$（控制周期） & 0.5 s & 邻域半径 & 500 m \\
        最大步长 & 200步（100 s） & 训练轮次 & 1000 \\
        目标速度 & $70\pm10$ km/h & $\gamma,\ \epsilon$ & 0.99, 0.2 \\
        速度范围 & $[60, 140]$ km/h & 初始学习率$\eta$ & $5\times10^{-4}$ \\
        最大加速度 & 3.0 m/s$^2$ & 最大减速度 & 4.0 m/s$^2$ \\
        风扰强度 & 0.1 m/s & 噪声标准差$\sigma$ & 0.5 m \\
        \hline
    \end{tabular}
\end{table}

\subsubsection{基线方法}

选取8种覆盖"规则调度—单/独立RL—CTDE-MARL—工程安全方法"四个层级的基线。规则调度包括轮询（RR）和先到先服务（FCFS），代表无学习能力的最朴素策略，用来衡量学习方法的增益下界；单/独立RL包括集中式DQN和独立A2C；CTDE-MARL包括MADDPG、COMA、QMIX和MAPPO，其中MAPPO是最强基线；工程方法为ORCA\cite{ref4}（pyrvo实现，时间前瞻$\tau=5.0$ s，安全半径与$d_{\text{safety}}$一致，首选速度$70\pm10$ km/h指向中心）。ORCA的前身VO在$N=21$下被ORCA一致主导，主实验不再单独报告。

所有RL基线与本文方法共享相同的环境与奖励函数，保证比较公平；规则与工程方法直接施加其固有策略。MPC和CBF未纳入比较，原因有二：MPC在21架飞行器场景下实时性难以满足0.5 s控制周期；CBF在高密度下安全集合交集容易退化。二者与MARL的混合架构是重要的未来方向。

\subsubsection{评价指标}

指标分"表层"和"深层"两个层级。表层指标刻画整体运行性能，包括累计奖励（CR，整个Episode的折扣累计奖励）、总冲突次数（TC，整个Episode所有时刻所有飞行器对之间的冲突数量累计）、平均延误（AD，相对理想飞行时间的平均延迟）。深层指标刻画安全性能：通过率（PCR，成功通过交叉中心的飞行器占比）、无冲突完成率（CFCR，全程零冲突且全部通过的Episode占比）、最小间隔违反率（MSVR，发生$d_{\min}<d_{\text{safety}}$的时步占比）、平均冲突幅度（MCAE，单Episode内平均每次冲突的严重程度）。

这里需要重点区分PCR与CFCR：PCR只衡量飞行器是否到达，CFCR则要求飞行器在不发生任何间隔违反的前提下到达。若只报告PCR，可能系统性高估模型的安全性能——这正是本文引入CFCR的核心动机。

\subsection{训练收敛性分析}

训练过程中的累计奖励、总冲突次数与课程进阶情况如图\ref{fig:training}所示，整体呈明显的三阶段特征。

\begin{figure}[H]
    \centering
    \includegraphics[width=0.8\textwidth]{QQ图片20260504231524.jpg}
    \caption{训练收敛曲线。阴影表示$\pm1$标准差（$n=5$），红色虚线标记课程学习进阶点。}
    \label{fig:training}
\end{figure}

\textbf{阶段I（0--200轮，快速上升期）：}累计奖励从约$-2000$升到$+800$，总冲突次数从80降到25。此时课程学习处于初级阶段（$N=5\sim11$），场景宽松，模型快速掌握基本避让策略。

\textbf{阶段II（200--600轮，进阶波动期）：}课程学习触发$N:11\to19$的多次进阶。每次进阶后，新加入的飞行器让场景密度骤增，模型出现1--3轮的短暂回落（图中红色虚线标记的进阶点），随后迅速适应。这种"进阶—回落—再适应"的锯齿形曲线，是课程学习正常工作的标志。

\textbf{阶段III（600--1000轮，收敛稳定期）：}累计奖励稳定在$1286.7\pm18.2$，总冲突次数收敛到$18.2\pm2.1$，在$N=21$的高密度场景下达到稳定性能。

收敛曲线也是训练稳定性的直接体现：若曲线剧烈振荡或长期不收敛，通常意味着训练机制存在问题。本文的三阶段收敛特征说明课程学习有效缓解了稀疏奖励导致的训练不稳定，为最终性能奠定了基础。

\subsection{综合性能对比}

表\ref{tab:comparison}给出21架eVTOL场景下各方法的综合性能（$n=5$，均值$\pm$标准差）。

\begin{table}[H]
    \centering
    \caption{综合性能对比（21架eVTOL，$n=5$）}
    \label{tab:comparison}
    \begin{tabular}{|c|c|c|c|c|}
        \hline
        \textbf{方法} & \textbf{CR$\uparrow$} & \textbf{TC$\downarrow$} & \textbf{AD(s)$\downarrow$} & \textbf{PCR(\%)$\uparrow$} \\
        \hline
        RR & $-1256.3\pm89.2$ & $42.6\pm5.3$ & $1.58\pm0.12$ & $89.2\pm3.5$ \\
        FCFS & $-1187.5\pm76.4$ & $38.9\pm4.7$ & $1.42\pm0.10$ & $91.5\pm2.8$ \\
        DQN & $-892.5\pm65.3$ & $28.3\pm3.9$ & $1.21\pm0.09$ & $94.7\pm2.1$ \\
        独立A2C & $-657.8\pm52.6$ & $19.7\pm2.8$ & $0.93\pm0.07$ & $96.3\pm1.7$ \\
        ORCA & $-198.7\pm35.8$ & $22.8\pm2.9$ & $0.76\pm0.06$ & $96.0\pm1.6$ \\
        MADDPG & $-215.6\pm41.2$ & $25.4\pm3.2$ & $0.87\pm0.06$ & $95.8\pm1.9$ \\
        COMA & $187.3\pm35.7$ & $22.1\pm2.9$ & $0.79\pm0.06$ & $97.1\pm1.5$ \\
        QMIX & $426.5\pm28.9$ & $23.1\pm2.7$ & $0.76\pm0.05$ & $97.5\pm1.3$ \\
        MAPPO & $892.4\pm23.5$ & $20.3\pm2.4$ & $0.65\pm0.04$ & $98.2\pm1.1$ \\
        \textbf{本文} & $\mathbf{1286.7\pm18.2}$ & $\mathbf{18.2\pm2.1}$ & $\mathbf{0.57\pm0.03}$ & $\mathbf{99.1\pm0.8}$ \\
        \hline
    \end{tabular}
\end{table}

从表中可以归纳出四条结论。第一，学习方法对规则方法呈系统性优势：RR的TC高达42.6、AD超过1.4 s、奖励为负，基本处于勉强运行的状态；本文相对RR将TC降低57.2\%、AD缩小63.9\%。第二，CTDE-MARL全面优于单智能体方法（DQN、独立A2C），再次验证集中训练对协作的重要性。第三，本文方法在四项指标上全部最优。第四，ORCA（TC=22.8）明显强于规则方法，但它的CFCR只有约50\%（见表\ref{tab:safety}），低于本文的68.4\%——ORCA缺乏多步轨迹预测能力，高密度下速度空间的可行解容易退化。

关于本文与ORCA的TC差异（18.2 vs.\ 22.8，降幅20.2\%）需要谨慎解读：以$n=5$和标准差估算，这个差异的统计显著性可能处于边界状态，系统性的比较要等$n\geq20$的样本。

\subsection{安全指标专项分析}

表层指标与深层安全指标的对比见表\ref{tab:safety}。

\begin{table}[H]
    \centering
    \caption{安全指标对比（$n=5$）}
    \label{tab:safety}
    \begin{tabular}{|c|c|c|c|c|}
        \hline
        \textbf{方法} & \textbf{CFCR(\%)$\uparrow$} & \textbf{MCAE$\downarrow$} & \textbf{MSVR(\%)$\downarrow$} & \textbf{PCR(\%)$\uparrow$} \\
        \hline
        RR & $12.3\pm4.1$ & $2.03\pm0.25$ & $8.7\pm1.2$ & $89.2\pm3.5$ \\
        FCFS & $18.7\pm3.6$ & $1.85\pm0.22$ & $7.4\pm1.0$ & $91.5\pm2.8$ \\
        ORCA & $50.2\pm3.0$ & $1.09\pm0.14$ & $3.8\pm0.6$ & $96.0\pm1.6$ \\
        DQN & $37.5\pm3.2$ & $1.35\pm0.19$ & $5.1\pm0.8$ & $94.7\pm2.1$ \\
        独立A2C & $48.2\pm2.8$ & $0.94\pm0.13$ & $3.6\pm0.6$ & $96.3\pm1.7$ \\
        MADDPG & $42.6\pm3.0$ & $1.21\pm0.15$ & $4.2\pm0.7$ & $95.8\pm1.9$ \\
        COMA & $51.3\pm2.7$ & $1.05\pm0.14$ & $3.8\pm0.6$ & $97.1\pm1.5$ \\
        QMIX & $53.7\pm2.5$ & $1.10\pm0.13$ & $3.5\pm0.5$ & $97.5\pm1.3$ \\
        MAPPO & $57.8\pm2.2$ & $0.97\pm0.11$ & $3.0\pm0.4$ & $98.2\pm1.1$ \\
        \textbf{本文} & $\mathbf{68.4\pm1.8}$ & $\mathbf{0.87\pm0.10}$ & $\mathbf{2.4\pm0.3}$ & $\mathbf{99.1\pm0.8}$ \\
        \hline
    \end{tabular}
\end{table}

\textbf{PCR与CFCR之间约30个百分点的系统性差距，是本文最重要的实证发现。}以本文方法为例，PCR高达99.1\%——几乎全部飞行器都完成了穿越——但CFCR只有68.4\%，意味着约三成Episode存在至少一次间隔违反。这个现象在MAPPO上更明显（PCR 98.2\% vs.\ CFCR 57.8\%，差约40个百分点），说明它不是本文方法的特例，而是CTDE-MARL在安全攸关评估中的共性挑战。

这个发现的方法论意义在于：如果安全评估只报PCR，会系统性高估模型的安全性能。CFCR把"只要不出事故"的宽松标准，提升为"全程零风险接触"的严格标准，更贴合低空交通安全的实际诉求。

\subsection{统计显著性检验}

为确认本文方法与最强基线MAPPO的差异不是随机波动，用Welch $t$检验\cite{ref42}（双侧，显著性水平$\alpha=0.05$，$n_1=n_2=5$）做统计检验。Welch $t$检验适用于两组样本方差可能不等的独立样本均值比较。除$p$值外，还报告Cohen's $d$效应量\cite{ref43}与统计功效。效应量以标准差为单位刻画两组均值差距的大小，是量纲无关的指标，一般$d>0.8$即视为大效应；统计功效指当真实差异确实存在时，检验能够正确发现该差异的概率，功效越高，越不容易漏报真实效应。结果见表\ref{tab:ttest}。

\begin{table}[H]
    \centering
    \caption{本文 vs.\ MAPPO统计检验与功效分析}
    \label{tab:ttest}
    \begin{tabular}{|c|c|c|c|c|c|c|}
        \hline
        \textbf{指标} & \textbf{均值差} & \textbf{$t$(df=8)} & \textbf{$p$} & \textbf{$d$} & \textbf{Power} & \textbf{结论} \\
        \hline
        TC & $-2.1$ & 1.47 & 0.180 & 0.93 & 38\% & 不显著（$n\geq20$需） \\
        CFCR & $+10.6\%$ & 8.34 & $<0.001$ & 5.27 & $>99\%$ & \textbf{高度显著} \\
        MCAE & $-0.10$ & 1.50 & 0.172 & 0.95 & 39\% & 不显著（$n\geq20$需） \\
        MSVR & $-0.6\%$ & 2.68 & 0.029 & 1.69 & 70\% & \textbf{显著} \\
        \hline
    \end{tabular}
\end{table}

\textbf{核心统计结论：}本文在CFCR（$p<0.001$，效力$>99\%$）和MSVR（$p=0.029$，效力70\%）两项安全指标上取得了统计可信的改善，这是本文的核心统计贡献。CFCR的效应量$d=5.27$远超"大效应"阈值0.8，说明$+10.6\%$的提升不仅在统计上显著，工程意义上也相当可观。

TC和MCAE的差异（$d=0.93$与0.95）是大效应量，但在$n=5$下统计效力只有约38\%，未达显著（$p=0.180$与0.172）。这里的$p=0.180$要恰当解读为"在现有样本量下没检测到显著差异"，而不是"确认无差异"——大效应量恰恰暗示真实效应很可能存在。据功效分析，需$n\geq20$才让检验达到80\%效力。这个区分对科学解读至关重要：既不能把不显著夸大成"确认无差异"，也不能忽视它的方向性证据价值。

\subsection{消融实验}

为厘清各模块的独立贡献，将完整模型与移除单一模块的变体对比，结果见表\ref{tab:ablation}。

\begin{table}[H]
    \centering
    \caption{消融实验（$n=5$）}
    \label{tab:ablation}
    \resizebox{0.92\textwidth}{!}{%
    \begin{tabular}{|c|c|c|c|c|}
        \hline
        \textbf{配置} & \textbf{CR$\uparrow$} & \textbf{TC$\downarrow$} & \textbf{AD(s)$\downarrow$} & \textbf{PCR(\%)$\uparrow$} \\
        \hline
        完整模型 & $\mathbf{1286.7\pm18.2}$ & $\mathbf{18.2\pm2.1}$ & $\mathbf{0.57\pm0.03}$ & $\mathbf{99.1\pm0.8}$ \\
        \hline
        移除DyGAT（标准GAT，无门控） & $215.3\pm32.6$ & $49.8\pm4.5$ & $0.79\pm0.05$ & $94.3\pm1.9$ \\
        \hline
        移除课程学习（直接$N=21$） & $587.2\pm27.4$ & $28.4\pm3.1$ & $0.68\pm0.04$ & $97.2\pm1.4$ \\
        \hline
        移除冲突损失加权（$\lambda\equiv1.0$） & $763.5\pm24.1$ & $24.6\pm2.8$ & $0.65\pm0.04$ & $97.8\pm1.2$ \\
        \hline
        同时移除课程学习与冲突加权 & $128.6\pm38.5$ & $35.7\pm3.9$ & $0.82\pm0.06$ & $95.6\pm1.7$ \\
        \hline
    \end{tabular}%
    }
\end{table}

消融结果传递了三条信息。第一，\textbf{DyGAT贡献最大}：去掉它（退化为无门控的标准GAT）后，TC从18.2上升到49.8（+173.6\%），CR从1286.7骤降到215.3，直接验证了乘法门控的核心价值。第二，课程学习与冲突加权的独立贡献：单独去掉课程学习，TC增加56.0\%；单独去掉冲突加权，TC增加35.2\%。第三，二者的联合效应：同时去掉时TC增加96.2\%，接近二者独立效应的加和（91.2\%），说明两个机制作用基本独立、可以放心叠加。

\subsection{密度鲁棒性分析}

前文实验都在$N=21$下进行。为检验密度鲁棒性，把飞行器数量从5逐步加到21，观察各方法TC的变化趋势（图\ref{fig:density}）。三条线的增长斜率差异显著：本文约0.75次/架，MAPPO约1.15次/架，RR高达约2.4次/架。

\begin{figure}[H]
    \centering
    \includegraphics[width=0.65\textwidth]{QQ图片20260524171218.png}
    \caption{不同密度下冲突次数对比。误差条表示$\pm1$标准差（$n=5$）。}
    \label{fig:density}
\end{figure}

斜率对比的工程含义很直观：密度每增加一架，本文只多约0.75次冲突，而规则方法要多2.4次。低密度时差距还不明显，一旦逼近真实运行场景的极限密度，差距会线性放大。DyGAT在高密度下的优势可归因于门控机制：邻居数量增多后，全连接图上的注意力信息急剧膨胀，无门控的方法容易被大量弱相关邻居稀释有效信息，而乘法门控将注意力聚焦于真正构成威胁的少数飞行器对，使模型在高密度下仍能保持对关键交互的聚焦。

\subsection{参数敏感性分析}

测试了三个关键超参数：基本安全间隔$d_{\text{base}}$、风扰强度与传感器噪声（图\ref{fig:sensitivity}）。

\begin{figure}[H]
    \centering
    \includegraphics[width=0.8\textwidth]{QQ图片20260524162300.png}
    \caption{超参数敏感性分析。红色虚线标记默认参数取值。误差条表示$\pm1$标准差。}
    \label{fig:sensitivity}
\end{figure}

\textbf{安全间隔-效率权衡：}$d_{\text{base}}$增大使TC降34.7\%（更安全），但AD升21.2\%（更延误）。这符合物理直觉：更大的间隔意味着更大的缓冲，但也限制了交通流的紧凑性。默认50 m在二者之间取了平衡。

\textbf{风扰的破坏性：}风扰强度增大使TC上升约28\%。风扰破坏航迹精确性，迫使模型频繁修正速度航向，冲突概率上升。

\textbf{噪声的有限影响：}传感器噪声增大使TC上升约15\%，三者中最小。原因在于LSTM对时序信息的平滑作用——历史序列的聚合天然过滤单步噪声，体现了时序建模的鲁棒性价值。

敏感性分析的意义在于确认性能的稳健边界：参数发生±20\%扰动时性能不会剧烈退化，说明性能并非依赖过拟合特定参数组合。

\subsection{推理效率分析}

低空管控对实时性要求苛刻：控制周期0.5 s，单步决策时间必须远小于它。各方法在不同$N$下的单步推理时间见图\ref{fig:computation}。

\begin{figure}[H]
    \centering
    \includegraphics[width=0.8\textwidth]{QQ图片20260524163903.jpg}
    \caption{不同飞行器数量$N$下单步推理时间对比（RTX 4060 GPU）。}
    \label{fig:computation}
\end{figure}

本文方法在$N=21$时单步推理4.5 ms，相对0.5 s的控制周期有约111倍的实时性裕度，即便计入通信与调度开销也仍有充足余量。相比其他MARL基线，本文推理开销处于同一量级，且随$N$近似线性增长。值得一提的是乘法门控的计算开销极低：它只涉及物理量的指数运算、无需梯度回传，在$N=21$时单步仅增加约0.02 ms，几乎可以忽略——物理先验注入在工程上是一种零成本的安全增强手段。

%==========================================================================
% 6 讨论与结论
%==========================================================================
\section{讨论与结论}

\subsection{与MS-GRL的比较}

Li等\cite{ref10}提出的MS-GRL是图增强MARL在低空无人机冲突消解领域的代表性工作，与本文具有较强的可比性，三个层面的差异值得逐一说明。算法上，MS-GRL用Double DQN（离散动作空间），本文用PPO（连续加速度）；连续动作能输出更平滑的加速度指令，避免离散动作导致的航迹抖动。图结构上，MS-GRL用固定距离阈值建图，凡在阈值半径内一律当邻居；本文用乘法门控的动态边权重，边权重随物理状态连续变化。场景规模上，MS-GRL处理100架UAV的编队冲突，本文处理21架eVTOL的交叉路口，前者是"大规模稀疏冲突"，后者是"小规模高密度冲突"，挑战重心不同。

门控的优势可以落到具体数字上：两对距离都是200 m的飞行器，一对同向等速（$w_{ij}\approx0.82$），一对相向高速（$w_{ij}\approx0.015$），注意力权重相差约55倍。固定阈值GAT把两者一视同仁，完全捕捉不到这个差异；DyGAT却实现了对"同样距离、不同风险"飞行器对的自适应区分。需要指出，MS-GRL与本文在路线上互补，把前者的多尺度图框架与本文的物理门控思想结合，是未来值得尝试的方向。

\subsection{核心发现与机制解读}

\textbf{CFCR揭示的安全评估盲区。}PCR（99.1\%）与CFCR（68.4\%）之间约30个百分点的差距，在MAPPO上更大（约40个百分点），说明只报PCR会系统性高估CTDE-MARL的安全性能。这个发现不限于本文方法，而是指出了该领域评估方法论的一个普遍盲区。

\textbf{乘法门控的贡献结构。}DyGAT整体替换让TC增加173.6\%，但其内部三个组件的独立贡献尚未分解。尤其是第4.3.3节提出的假设——$w_{ij}$在学习注意力产生偏差时提供必要约束——引出一个待验证的问题：如果只用$w_{ij}$做注意力（完全不学$e_{ij}$），性能会接近完整DyGAT吗？回答它需要"仅$w_{ij}$ / 仅$e_{ij}$ / $e_{ij}\odot w_{ij}$"的三路组件级消融。

\textbf{统计证据的差异化图景。}本文的核心统计贡献定位于CFCR（$p<0.001$，$d=5.27$）和MSVR（$p=0.029$，$d=1.69$）的显著改善，这两项在安全攸关评估里工程意义更直接。TC的差异（$p=0.180$，$d=0.93$）应视为方向性证据——大效应量说明真实效应很可能存在，但$n=5$的统计效力不足以确认。把论述重心从TC移到CFCR/MSVR，是对统计证据强度的诚实反映。

\subsection{研究局限性}

本文工作仍存在若干局限，需要客观认识。

\textbf{（1）二维场景限制。}当前是二维单交叉路口，尚未涉及三维高度层、多路口网络和起降场约束。真实低空交通的复杂度远超这个抽象场景，三维扩展是首要的后续工作。

\textbf{（2）统计样本量。}$n=5$对TC这类中等效应量指标检验力不足（38\%），$n\geq20$的重复实验是必要的验证步骤。

\textbf{（3）风扰与噪声模型。}均匀风扰和高斯白噪声都是时间独立的简化模型。真实风场用Dryden/Kaimal谱模型描述，具有空间相关性和时间连续性，对航迹的影响更复杂。

\textbf{（4）通信假设。}CTDE训练假设全局状态无延迟获取，实际UTM部署里通信有延迟（50--200 ms）和丢包，可能削弱集中训练的效果。

\textbf{（5）基线覆盖。}MPC和CBF尚未纳入比较；ORCA做了场景适配调优，但没对所有基线做系统性超参数搜索，基线的性能上限可能没被充分挖掘。

\textbf{（6）eVTOL动力学参数。}加速度取值（3.0--4.0 m/s$^2$）是参考学术文献\cite{ref25}在缺乏厂商精确数据下的估计值，拿到厂商数据后需做敏感性验证。

\textbf{（7）门控超参数。}三个指数核的衰减常数（1000 m、20 km/h、$\pi/4$ rad）为手工预设，未系统调优，可能限制模型在不同空域尺度下的泛化。

\subsection{结论}

本文针对城市低空十字交叉路口的eVTOL协同管控问题，提出了一种融合乘法门控注意力掩码的耦合多智能体强化学习模型。核心创新是DyGAT中基于物理量的乘法门控机制——用距离、速度差和航向差编码出软性注意力掩码，实现飞行器间时空耦合关系的细粒度建模。配套的速度自适应安全距离、分层奖励、课程学习与冲突损失加权，共同解决了稀疏冲突样本下的训练困难。

在21架eVTOL的二维交叉路口仿真中，本文方法在CFCR（$68.4\%$ vs.\ MAPPO $57.8\%$，$p<0.001$，$d=5.27$）和MSVR（$2.4\%$ vs.\ $3.0\%$，$p=0.029$，$d=1.69$）上取得统计显著的改善；TC的差异（$18.2$ vs.\ $20.3$，降幅10.3\%）有大效应量（$d=0.93$）但受$n=5$限制未达显著，有待$n\geq20$的重复实验确认。消融实验表明DyGAT是贡献最大的模块。

后续工作主要包括：（1）三维多路口网络扩展；（2）$n\geq20$重复实验并补充vs.\ ORCA的完整统计比较；（3）集成Dryden/Kaimal标准风谱模型；（4）引入通信延迟与丢包约束；（5）DyGAT三路组件级消融（仅$w_{ij}$/仅$e_{ij}$/$e_{ij}\odot w_{ij}$）；（6）与MPC/CBF的系统基准对比。

\textbf{代码与可复现性声明：}实验代码（含AirTrafficEnv仿真环境、DyGAT+GAT-Mixer模型实现、ORCA基线适配、训练与评估脚本）将在论文接收后于GitHub开源（匿名链接在camera-ready版本提供）。复现本文结果约需1000轮$\times$400 episodes$\times$4.5 ms$\times$1000步$\approx$5小时（RTX 4060 GPU）每seed，5 seeds约25 GPU小时，$n=20$约需100 GPU小时。

%==========================================================================
% 参考文献
%==========================================================================
\begin{thebibliography}{99}

\bibitem{ref1} Federal Aviation Administration. Urban air mobility (UAM): concept of operations v2.0[R]. Washington, DC: FAA, 2023.

\bibitem{ref2} MENON P K, SWERIDUK G D, BILIMORIA K D. New approach to modeling, analysis, and control of air traffic flow[J]. Journal of Guidance, Control, and Dynamics, 2004, 27(5): 737--744.

\bibitem{ref3} RICHARDS A, HOW J P. Aircraft trajectory planning with collision avoidance using mixed integer linear programming[C]//Proceedings of the 2002 American Control Conference. Piscataway: IEEE, 2002: 1936--1941.

\bibitem{ref4} VAN DEN BERG J, GUY S J, LIN M, et al. Reciprocal n-body collision avoidance[C]//Robotics Research: The 14th International Symposium ISRR (Springer Tracts in Advanced Robotics, Vol.\ 70). Berlin: Springer, 2011: 3--19.

\bibitem{ref5} AMES A D, COOGAN S, EGERSTEDT M, et al. Control barrier functions: theory and applications[C]//2019 18th European Control Conference (ECC). Piscataway: IEEE, 2019: 3420--3431.

\bibitem{ref6} AGOGINO A, TUMER K. Regulating air traffic flow with coupled agents[C]//Proceedings of the 7th International Conference on Autonomous Agents and Multiagent Systems (AAMAS 2008). Richland: IFAAMAS, 2008: 535--542.

\bibitem{ref7} AGOGINO A K, TUMER K. A multiagent approach to managing air traffic flow[J]. Autonomous Agents and Multi-Agent Systems, 2012, 24(1): 1--25.

\bibitem{ref8} BRITTAIN M, WEI P. Autonomous separation assurance in a high-density en route sector: a deep multi-agent reinforcement learning approach[C]//2019 IEEE Intelligent Transportation Systems Conference (ITSC). Piscataway: IEEE, 2019: 3256--3262.

\bibitem{ref9} GHOSH S, LAGUNA S, LIM S H, et al. A deep ensemble method for multi-agent reinforcement learning: a case study on air traffic control[C]//Proceedings of the 31st International Conference on Automated Planning and Scheduling (ICAPS 2021). Palo Alto: AAAI Press, 2021: 468--476.

\bibitem{ref10} LI Y, LI J, WANG J, ZHANG X, DING H, DU W. Multiscale-graph-enhanced reinforcement learning for conflict resolution in dense UAV networks[J]. IEEE Internet of Things Journal, 2025, 12(21): 44290--44303.

\bibitem{ref11} SPATHARIS C, BASTAS A, KRAVARIS T, et al. Hierarchical multiagent reinforcement learning schemes for air traffic management[J]. Neural Computing and Applications, 2021, 33(14): 8649--8671.

\bibitem{ref12} ASLANI M, MESGARI M S, WIERING M. Adaptive traffic signal control with actor-critic methods in a real-world traffic network with different traffic disruption events[J]. Transportation Research Part C: Emerging Technologies, 2017, 85: 732--752.

\bibitem{ref13} 翟子洋, 郝茹茹, 董世浩. 大规模智慧交通信号控制中的强化学习和深度强化学习方法综述[J]. 计算机应用研究, 2024, 41(6): 1618--1627.

\bibitem{ref14} TUMER K, AGOGINO A. Distributed agent-based air traffic flow management[C]//Proceedings of the 6th International Joint Conference on Autonomous Agents and Multiagent Systems (AAMAS 2007). Honolulu: ACM, 2007: 330--337.

\bibitem{ref15} 王迎. 城市交通信号控制深度强化学习算法研究[D]. 济南: 山东交通学院, 2024.

\bibitem{ref16} 赵晓华, 石建军, 李振龙, 赵国勇. 基于Q-learning和BP神经元网络的交叉口信号灯控制[J]. 公路交通科技, 2007, 24(7): 99--102.

\bibitem{ref17} 赖建辉. 基于深度强化学习的交通控制优化方法研究及实现[D]. 杭州: 浙江工业大学, 2020.

\bibitem{ref18} 张新武. 基于深度强化学习算法的交通信号控制优化研究[D]. 太原: 太原科技大学, 2024.

\bibitem{ref19} 承向军, 常歆识, 杨肇夏. 基于Q-学习的交通信号控制方法[J]. 系统工程理论与实践, 2006(8): 136--140.

\bibitem{ref20} 卢守峰, 张术, 刘喜敏. 平均排队长度差最小的单交叉口在线Q学习模型[J]. 公路交通科技, 2014, 31(11): 116--122.

\bibitem{ref21} 江波, 李彦冬, 刘威陇, 等. 应用深度Q学习的机场道口协同智能调速方法[J]. 航空计算技术, 2022, 52(1): 26--30.

\bibitem{ref22} 张春阳, 席雅琪. 人工智能在城市交通控制中的应用综述[J]. 信息系统工程, 2024(07): 33--36.

\bibitem{ref23} GHAVAMZADEH M, MAHADEVAN S, MAKAR R. Hierarchical multi-agent reinforcement learning[J]. Autonomous Agents and Multi-Agent Systems, 2006, 13(2): 197--229.

\bibitem{ref24} TUMER K, AGOGINO A. Improving air traffic management with a learning multiagent system[J]. IEEE Intelligent Systems, 2009, 24(1): 18--21.

\bibitem{ref25} BACCHINI A, CESTINO E. Electric VTOL configurations comparison[J]. Aerospace, 2019, 6(3): 26.

\bibitem{ref26} U.S. Department of Defense. Flying qualities of piloted airplanes: MIL-F-8785C[S]. Washington, DC: U.S. Department of Defense, 1980.

\bibitem{ref27} LIANG X, MA Y, FENG Y, et al. PTR-PPO: proximal policy optimization with prioritized trajectory replay[EB/OL]. arXiv:2112.03798, 2021.

\bibitem{ref28} CONG W, ZHANG S, CHEN J, et al. Do we really need complicated model architectures for temporal networks?[C]//Proceedings of the 11th International Conference on Learning Representations (ICLR 2023). Kigali: ICLR, 2023.

\bibitem{ref29} XIE Y B, GARDI A, SABATINI R. Reinforcement learning-based flow management techniques for urban air mobility and dense low-altitude air traffic operations[C]//2021 IEEE/AIAA 40th Digital Avionics Systems Conference (DASC). Piscataway: IEEE, 2021: 1--10.

\bibitem{ref30} SCHULMAN J, WOLSKI F, DHARIWAL P, et al. Proximal policy optimization algorithms[EB/OL]. arXiv:1707.06347, 2017.

\bibitem{ref31} VELI\v{C}KOVI\'C P, CUCURULL G, CASANOVA A, et al. Graph attention networks[C]//Proceedings of the 6th International Conference on Learning Representations (ICLR 2018). Vancouver: ICLR, 2018.

\bibitem{ref32} KIPF T N, WELLING M. Semi-supervised classification with graph convolutional networks[C]//Proceedings of the 5th International Conference on Learning Representations (ICLR 2017). Toulon: ICLR, 2017.

\bibitem{ref33} RASHID T, SAMVELYAN M, DE WITT C S, et al. QMIX: monotonic value function factorisation for deep multi-agent reinforcement learning[C]//Proceedings of the 35th International Conference on Machine Learning (ICML 2018). Stockholmsm\"{a}ssan: PMLR, 2018: 4295--4304.

\bibitem{ref34} LOWE R, WU Y, TAMAR A, et al. Multi-agent actor-critic for mixed cooperative-competitive environments[C]//Advances in Neural Information Processing Systems 30 (NeurIPS 2017). Long Beach: Curran Associates, 2017: 6379--6390.

\bibitem{ref35} FOERSTER J, FARQUHAR G, AFOURAS T, et al. Counterfactual multi-agent policy gradients[C]//Proceedings of the 32nd AAAI Conference on Artificial Intelligence (AAAI 2018). New Orleans: AAAI Press, 2018: 2974--2982.

\bibitem{ref36} YU C, VELU A, VINITSKY E, et al. The surprising effectiveness of PPO in cooperative multi-agent games[C]//Advances in Neural Information Processing Systems 35 (NeurIPS 2022). New Orleans: Curran Associates, 2022: 24611--24624.

\bibitem{ref37} HOCHREITER S, SCHMIDHUBER J. Long short-term memory[J]. Neural Computation, 1997, 9(8): 1735--1780.

\bibitem{ref38} SCHULMAN J, MORITZ P, LEVINE S, et al. High-dimensional continuous control using generalized advantage estimation[C]//Proceedings of the 4th International Conference on Learning Representations (ICLR 2016). San Juan: ICLR, 2016.

\bibitem{ref39} BENGIO Y, LOUGARD J, COLLOBERT R, et al. Curriculum learning[C]//Proceedings of the 26th International Conference on Machine Learning (ICML 2009). Montreal: ACM, 2009: 41--48.

\bibitem{ref40} MNIH V, KAVUKCUOGLU K, SILVER D, et al. Human-level control through deep reinforcement learning[J]. Nature, 2015, 518(7540): 529--533.

\bibitem{ref41} BERNSTEIN D S, GIVAN R, IMMERMAN N, et al. The complexity of decentralized control of Markov decision processes[J]. Mathematics of Operations Research, 2002, 27(4): 819--840.

\bibitem{ref42} WELCH B L. The generalization of Student's problem when several different population variances are involved[J]. Biometrika, 1947, 34(1--2): 28--35.

\bibitem{ref43} COHEN J. Statistical power analysis for the behavioral sciences[M]. 2nd ed. Hillsdale: Lawrence Erlbaum Associates, 1988.

\bibitem{ref44} International Electrotechnical Commission. Wind energy generation systems -- Part 1: design requirements: IEC 61400-1[S]. 4th ed. Geneva: IEC, 2019.

\bibitem{ref45} COHEN A P, SHAHEEN S A, FARRAR E M. Urban air mobility: history, ecosystem, market potential, and challenges[J]. IEEE Transactions on Intelligent Transportation Systems, 2021, 22(9): 6074--6087.

\end{thebibliography}

\end{document}
